
\section{Knowledge Graph and Database}
\label{sec:knowledge-graph-and-database}

This section describes the two data stores used in MLguide. The first 
is the ontology-based ML knowledge graph stored in GraphDB, which is based on ML 
Ontology and extended for use in MLguide. The second is the PostgreSQL 
database, which stores article abstracts and vector embeddings used 
for semantic retrieval.

\subsection{ML Ontology Extension}
\label{sec:ml-ontology-extension}

ML Ontology was selected as the knowledge base for MLguide because it 
provides the broadest coverage of ML concepts among the ontologies 
reviewed in Section~\ref{sec:ml-ontology-structure-and-content}, and because its 
structure directly supports the kind of structured method retrieval 
that MLguide requires. It is open source and published under an MIT 
licence \parencite{humm2026industry_ready_machine_learning_ontology}, 
which allows for use and modification as needed.
The limitations identified in 
Section~\ref{sec:ml-ontology-structure-and-content}, including the absence of 
reinforcement learning instances and incomplete coverage of 
\texttt{ML\_approach}, required extension before the ontology could 
be used in MLguide. Some extensions were also made to support better
knowledge structure and specific MLguide functionality. 
The extension was carried out gradually across the system
versions. This section describes the extension as a whole.


Three new classes were added under a separate namespace 
(\texttt{mla:}), short for \textit{machine learning articles}, 
to make clear that they are not part of the original ML Ontology. 
Table~\ref{tab:mla-classes} describes these classes.

\begin{table}[H]
\centering
\caption[New ontology classes]{New classes added under the \texttt{mla:} 
namespace.}
\label{tab:mla-classes}
\begin{tabularx}{\textwidth}{lX}
\toprule
\textbf{Class} & \textbf{Description} \\
\midrule
\texttt{mla:Article} & Represents a research article from the 
literature dataset. Linked to one or more \texttt{mla:Method} 
instances that are mentioned in the article. \\
\midrule
\texttt{mla:Method} & Represents an ML method as extracted from 
article abstracts. Linked to a corresponding 
\texttt{ML\_approach} instance via \texttt{skos:exactMatch}. \\
\midrule
\texttt{mla:ML\_task\_family} & Groups related \texttt{ML\_task} 
instances under a common family label, helps selecting notebook 
templates during notebook generation. \\
\bottomrule
\end{tabularx}
\end{table}

\texttt{mla:Method} and \texttt{ML\_approach} represent the same 
concept but serve different roles. The reason for keeping them 
separate is practical: \texttt{ML\_approach} instances carry 
relations to other ML concepts in the ontology, such as which 
tasks a method can be used for and which learning area it belongs 
to. Adding article relations directly to these instances would mix 
two different kinds of knowledge on the same node, making the 
ontology harder to navigate during development. Instead, article 
relations are kept on \texttt{mla:Method}, and the two classes 
are connected via \texttt{skos:exactMatch}. If needed, this 
separation could be removed by transferring the article relations 
directly to the corresponding \texttt{ML\_approach} instances.

Figure~\ref{fig:ontology-uml-with-extension} shows the class structure 
of the extended ontology. 

\begin{figure}[H]
    \centering
    \includegraphics[width=\textwidth]{Figures/06-implementation/ontology-uml-with-extension.pdf}
    \caption[ML Ontology UML diagram]{Unified Modeling Language (UML) class diagram of the ML Ontology, including added classes and properties for MLguide.}
    \label{fig:ontology-uml-with-extension}
\end{figure}

The upper part shows the classes from the 
original ML Ontology that are directly connected to the three new 
\texttt{mla:} classes: \texttt{ML\_approach}, \texttt{ML\_task}, and 
\texttt{ML\_area}. The lower part shows the new \texttt{mla:} classes 
and how they connect to these. \texttt{mla:Method} is linked to 
\texttt{ML\_approach} via \texttt{skos:exactMatch}, and to 
\texttt{mla:Article} via \texttt{mentionsMethod}. 
\texttt{mla:ML\_task\_family} is linked to \texttt{ML\_task} via 
\texttt{has\_family}, and \texttt{mla:Article} is linked to 
\texttt{ML\_area} via \texttt{has\_ml\_area}.

Figure~\ref{fig:ontology-article-instance} shows a concrete example 
of how the extended ontology represents an article and its relations. 

\begin{figure}[H]
    \centering
    \includegraphics[width=\textwidth]{Figures/06-implementation/ontology-article-instance.pdf}
    \caption[Article instance example]{Example of an Article instance in the ML Ontology.}
    \label{fig:ontology-article-instance}
\end{figure}

The \texttt{mla:Article} instance represents a research article 
identified by its Digital Object Identifier (DOI) and title. It is linked to 
\texttt{supervised\_learning} via \texttt{has\_ml\_area}, and to 
\texttt{method\_linear\_regression}, an \texttt{mla:Method} instance, 
via \texttt{mentionsMethod}. The method instance is in turn linked to 
the \texttt{ML\_approach} instance \texttt{linear\_regression} via 
\texttt{skos:exactMatch}, which carries the ontology relations: 
\texttt{used\_for} links it to the \texttt{ML\_task} 
\texttt{regression}, which belongs to \\ 
\texttt{supervised\_learning} 
via \texttt{belongs\_to}, and \texttt{has\_family} links it to the 
\texttt{mla:ML\_task\_family} instance \texttt{regression\_family}.

New \texttt{ML\_approach} instances were added to address the 
coverage gaps identified in Section~\ref{sec:ml-ontology-structure-and-content}. 
A first set was derived from the method extraction results of the 
pre-master project \parencite{bergsto2026mlselection}, where ML 
methods were identified from article abstracts in the literature 
dataset. Methods that appeared in the dataset but had no 
corresponding \texttt{ML\_approach} instance in ML Ontology were 
added. Reinforcement learning methods and tasks were added based 
on established taxonomies of RL algorithms 
\parencite{zhang2020taxonomy_reinforcement_learning_algorithms, 
geron2019hands_on_machine_learning}, following the same 
classification used in Section~\ref{sec:machine-learning-methods}. 
In total, 65 new instances were added, grouped by theme in 
Table~\ref{tab:new-ml-approaches}.

\begin{table}[H]
\centering
\caption[New ML\_approach instances]{New \texttt{ML\_approach} instances added to ML Ontology, grouped by theme.}
\label{tab:new-ml-approaches}
\begin{tabularx}{\textwidth}{lXr}
\toprule
\textbf{Theme} & \textbf{New instances} & \textbf{n} \\
\midrule
Reinforcement
  & deep reinforcement learning, Q-learning, REINFORCE, deep Q-network, policy gradient, proximal policy optimization, soft actor-critic, actor-critic, deep deterministic policy gradient, twin delayed DDPG, A3C & 11 \\
\midrule
Deep learning
  & deep neural network, attention mechanism, graph attention network, graph neural network, graph convolutional network, GraphSAGE, U-Net, variational autoencoder, normalizing flow, diffusion model, DeepAR & 11 \\
\midrule
Evolutionary
  & genetic algorithm, particle swarm optimization, simulated annealing, differential evolution, ant colony optimization, firefly algorithm, harmony search, evolutionary strategy, CMA-ES & 9 \\
\midrule
Anomaly detection
  & isolation forest, local outlier factor, one-class SVM, HBOS, Hotelling T\textsuperscript{2}, squared prediction error, minimum covariance determinant, elliptic envelope & 8 \\
\midrule
Time series
  & Kalman filter, temporal convolutional network, N-BEATS, seasonal ARIMA, state space model, VAR, Holt-Winters & 7 \\
\midrule
Explainability
  & explainable AI, SHAP, Granger causality, DoWhy, Causal Impact, causal forest & 6 \\
\midrule
Probabilistic
  & Bayesian method, Gaussian process regression, partial least squares, elastic net, robust regression & 5 \\
\midrule
Clustering
  & k-medoids, Gaussian mixture model clustering, t-SNE, UMAP & 4 \\
\midrule
Ensemble
  & gradient boosting, extra trees, bagging & 3 \\
\midrule
Other
  & naive Bayes & 1 \\
\midrule
\textbf{Total} & & \textbf{65} \\
\bottomrule
\end{tabularx}
\end{table}

For methods that form hierarchical families, \texttt{skos:broader} 
relations were added to preserve the inheritance structure of the 
ontology. For example, \texttt{deep\_q\_network} was linked as 
broader to \texttt{q\_learning}, and 
\texttt{proximal\_policy\_optimization} to \texttt{policy\_gradient}. 
After adding new instances, the inheritance rule described in 
Section~\ref{sec:ml-ontology-structure-and-content} was applied to propagate 
\texttt{used\_for} and \texttt{belongs\_to} relations to the new 
instances through their \texttt{skos:broader} hierarchies.

Where available, links to corresponding Wikidata entities were added 
to new instances via \texttt{rdfs:seeAlso}, following the same 
approach used in the original ontology 
\parencite{humm2026industry_ready_machine_learning_ontology}. 
Candidates were retrieved using the Wikidata search API and scored 
by string similarity between labels and whether the Wikidata 
description contained predefined ML-related terms. Results were 
reviewed manually before being added.

Table~\ref{tab:ml-approach-coverage-extended} shows the resulting 
coverage across all 187 \texttt{ML\_approach} instances after the 
extension and applying the inference rule. Coverage here means the 
share of \texttt{ML\_approach} instances that have a value for a 
given property. For reference, the original coverage across 122 
instances is shown in Table~\ref{tab:ml-approach-coverage}.

\begin{table}[H]
\centering
\caption[ML\_approach property coverage after extension]{Coverage 
of \texttt{ML\_approach} properties across 187 instances after 
ontology extension and inference rule application.}
\label{tab:ml-approach-coverage-extended}
\begin{tabular}{lr}
\toprule
\textbf{Property} & \textbf{Coverage} \\
\midrule
\texttt{skos:prefLabel}        & 187/187 (100\%) \\
\texttt{used\_for}            & 180/187 (96\%)  \\
\texttt{rdfs:comment}          & 130/187 (70\%)  \\
\texttt{performance}          & 95/187 (51\%)   \\
\texttt{rdfs:seeAlso}          & 90/187 (48\%)   \\
\texttt{skos:altLabel}         & 66/187 (35\%)   \\
\texttt{skos:broader}          & 57/187 (30\%)   \\
\texttt{possible\_if}         & 28/187 (15\%)   \\
\texttt{not\_recommended\_if} & 5/187 (3\%)     \\
\bottomrule
\end{tabular}

\par\vspace{0.5em}
{\footnotesize Note: \texttt{has\_dataset\_type}, which relates to 
\texttt{ML\_approach} indirectly through \texttt{ML\_task}, is 
populated for 29 of the 43 \texttt{ML\_task} instances that exist 
after the extension.}
\end{table}

Populating instances with properties requires manual work, and gathering 
the knowledge to do so takes effort. 
The knowledge used to update the knowledgegraph
was based on the established ML sources and textbooks also used in 
Section~\ref{sec:machine-learning}. When updating the knowledge, new instances were prioritised, since 
they had no properties at all. Reinforcement learning is one 
example, as it had no instances in the original ontology. Clear 
gaps found during testing of the recommendation process, or 
through user testing feedback, were also prioritised. Among the properties, \texttt{used\_for} and 
\texttt{performance} were prioritised, as these are central to the 
candidate retrieval and ranking described in 
Section~\ref{sec:ontology-candidate-retrieval}. As a result, 
coverage is improved but remains incomplete for several of the remaining 
properties, as Table~\ref{tab:ml-approach-coverage-extended} shows.

\texttt{mla:ML\_task\_family} instances were also added to group related 
\texttt{ML\_task} instances under a common label. This was motivated by
the large amount of similar tasks, especially for classification and regression. 
Grouping them into families provides 
a clearer structure. Task families also support notebook 
generation by providing a better structure for selecting the correct 
template for a given method and task combination, since many methods 
can be applied to more than one type of task. This is described in 
more detail in Section~\ref{sec:notebook-generation}.

Families were defined by inspecting the \texttt{skos:broader} 
hierarchy to identify which tasks were already grouped under common 
parent tasks, and also by grouping tasks of similar character where no 
such hierarchy existed. Each family was given a name that reflects 
the shared nature of its tasks.
Table~\ref{tab:ml-task-families} lists the task families and their 
assigned tasks.

\begin{table}[H]
\centering
\caption[ML\_task\_family instances]{ML\_task\_family instances and their assigned ML\_task instances.}
\label{tab:ml-task-families}
\begin{tabularx}{\textwidth}{>{\raggedright\arraybackslash}p{4cm}Xr}
\toprule
\textbf{ML\_task\_family} & \textbf{Assigned ML\_task instances} & \textbf{n} \\
\midrule
classification\_\allowbreak family
& classification, binary\_classification, multiclass\_classification,
    tabular\_classification, text\_classification, image\_classification,
    audio\_classification, time\_series\_classification, video\_classification,
    vertex\_classification & 10 \\
\midrule
regression\_\allowbreak family
& regression, tabular\_regression, text\_regression, image\_regression,
    audio\_regression, video\_regression & 6 \\
\midrule
clustering\_\allowbreak unsupervised\_\allowbreak discovery\_\allowbreak family
& clustering, tabular\_clustering, association\_rule\_learning,
    community\_detection, topic\_modeling & 5 \\
\midrule
detection\_\allowbreak localization\_\allowbreak family
& object\_detection, anomaly\_detection, action\_recognition & 3 \\
\midrule
feature\_\allowbreak engineering\_\allowbreak family
& feature\_extraction, feature\_reduction, feature\_selection & 3 \\
\midrule
structured\_\allowbreak prediction\_\allowbreak family
& image\_segmentation, named\_entity\_recognition, translation & 3 \\
\midrule
graph\_tasks\_\allowbreak family
& graph\_matching, link\_prediction, vertex\_classification & 3 \\
\midrule
nlp\_analysis\_\allowbreak family
& sentiment\_analysis, text\_analytics & 2 \\
\midrule
ranking\_\allowbreak recommendation\_\allowbreak similarity\_\allowbreak family
& ranking, recommendation, sentence\_similarity & 3 \\
\midrule
time\_series\_\allowbreak family
& time\_series\_analysis, time\_series\_forecasting, time\_series\_classification & 3 \\
\midrule
reinforcement\_\allowbreak learning\_\allowbreak family
& policy\_learning, value\_estimation, policy\_optimization & 3 \\
\midrule
tracking\_\allowbreak family
& video\_object\_tracking & 1 \\
\bottomrule
\end{tabularx}
\end{table}

\subsection{Database Structure}
\label{sec:database-structure}

Article data is stored across several tables in PostgreSQL, shown in 
Figure~\ref{fig:psql-er-diagram}. The \texttt{article\_abstracts} 
table stores the raw abstract text for each article, used during 
reranking. The \texttt{abstract\_embeddings} and 
\texttt{title\_embeddings} tables store 768-dimensional vector 
embeddings of the abstract and title respectively, used for semantic 
similarity search via pgvector \parencite{pgvector}. The 
\texttt{users} table stores a minimal user record created on first 
login. The \texttt{saved\_searches} table stores previously saved searches 
with problem descriptions, allowing users to reload them without 
re-entering their inputs. The \texttt{method\_feedback} table stores
user feedback on recommended methods, capturing all context related to 
the user's problem description and the specific method. The \texttt{skipped\_articles}
table is used in the process retrieving new articles for the knowledge graph and article database. 
It is further explained in Section~\ref{sec:article-retrieval-flowchart}.

\begin{figure}[H]
    \centering
    \includegraphics[width=\textwidth]{Figures/06-implementation/psql-er-diagram.pdf}
    \caption[PostgreSQL ER diagram]{Entity-relationship diagram of the PostgreSQL database used in MLguide.}
    \label{fig:psql-er-diagram}
\end{figure}